\documentclass[../../../include/open-logic-section]{subfiles}

\begin{document}

\olfileid{sth}{z}{power}
\olsection{في مجموعات القوى}

ونتبع ما تقدم بمسلّمة أخرى:

\begin{axiom}[مجموعات القوى]
لكل مجموعة~$\OLId{latin-looped}{A}$ توجد المجموعة~$\Pow{\OLId{latin-looped}{A}} = \Setabs{\OLId{latin-ordinary}{x}}{\OLId{latin-ordinary}{x} \subseteq \OLId{latin-looped}{A}}$.
\[
	\forall \OLId{latin-looped}{A} \exists \OLId{latin-looped}{P} \forall \OLId{latin-ordinary}{x}(\OLId{latin-ordinary}{x} \in \OLId{latin-looped}{P} \liff (\forall \OLId{latin-ordinary}{z} \in \OLId{latin-ordinary}{x})\OLId{latin-ordinary}{z} \in \OLId{latin-looped}{A})
\]
\end{axiom}

وتسويغها قريب. افترض أن~$\OLId{latin-looped}{A}$ تكونت في المرحلة~$\OLId{latin-looped}{S}$؛ فتكون عناصر~$\OLId{latin-looped}{A}$
كلها متاحة قبل~$\OLId{latin-looped}{S}$ بموجب~\stagesacc. وبحجة نظيرة لحجة الفصل تكون كل
مجموعة جزئية من~$\OLId{latin-looped}{A}$ قد تكونت عند المرحلة~$\OLId{latin-looped}{S}$، فتتاح كلها لتُجمع في
مجموعة واحدة عند أي مرحلة تلي~$\OLId{latin-looped}{S}$. ومثل هذه المرحلة موجودة، لأن~$\OLId{latin-looped}{S}$
ليست الأخيرة بموجب~\stagessucc. فثبت وجود~$\Pow{\OLId{latin-looped}{A}}$.

ومن النتائج الحسنة لمسلّمة مجموعات القوى:

\begin{prop}\label{thm:Products}
لكل مجموعتين~$\OLId{latin-looped}{A},\OLId{latin-looped}{B}$ يوجد حاصل ضربهما الديكارتي~$\OLId{latin-looped}{A} \times \OLId{latin-looped}{B}$.
\end{prop}

\begin{proof}
بمجموعات القوى و\olref[sth][z][pairs]{prop:pairsconsequences} توجد
$\Pow{\Pow{\OLId{latin-looped}{A} \cup \OLId{latin-looped}{B}}}$؛ ومن ثم يثبت بالفصل وجود المجموعة:
\[
	\OLId{latin-looped}{C} = \Setabs{\OLId{latin-ordinary}{z} \in \Pow{\Pow{\OLId{latin-looped}{A} \cup \OLId{latin-looped}{B}}}}{(\exists \OLId{latin-ordinary}{x} \in \OLId{latin-looped}{A})(\exists \OLId{latin-ordinary}{y} \in \OLId{latin-looped}{B}) \OLId{latin-ordinary}{z} = \tuple{\OLId{latin-ordinary}{x}, \OLId{latin-ordinary}{y}}}.
\]
ولكل~$\OLId{latin-ordinary}{x} \in \OLId{latin-looped}{A}$ و$\OLId{latin-ordinary}{y} \in \OLId{latin-looped}{B}$ توجد~$\tuple{\OLId{latin-ordinary}{x},\OLId{latin-ordinary}{y}}$ بموجب
\olref[sth][z][pairs]{prop:pairsconsequences}. ثم لما كان
$\OLId{latin-ordinary}{x},\OLId{latin-ordinary}{y} \in \OLId{latin-looped}{A} \cup \OLId{latin-looped}{B}$، تحقق~$\{\OLId{latin-ordinary}{x}\},\{\OLId{latin-ordinary}{x},\OLId{latin-ordinary}{y}\} \in \Pow{\OLId{latin-looped}{A} \cup \OLId{latin-looped}{B}}$، ومنه
$\tuple{\OLId{latin-ordinary}{x},\OLId{latin-ordinary}{y}} \in \Pow{\Pow{\OLId{latin-looped}{A} \cup \OLId{latin-looped}{B}}}$. فـ~$\OLId{latin-looped}{A} \times \OLId{latin-looped}{B} = \OLId{latin-looped}{C}$.
\end{proof}

وفي البرهان تعاون بين مسلّمة مجموعات القوى والفصل، وليس ذلك بعجيب؛ إذ
لولا الفصل لم تكن مجموعات القوى مبدأ شديد \emph{القوة}. فالفصل يعيّن
المجموعات الجزئية الموجودة من كل مجموعة، ومن ثم يعيّن «عظم» مجموعة
قواها.

\begin{prob}
أثبت، لكل مجموعتين~$\OLId{latin-looped}{A},\OLId{latin-looped}{B}$: (i) وجود مجموعة العلاقات كلها التي مجالها~$\OLId{latin-looped}{A}$
ومداها~$\OLId{latin-looped}{B}$؛ و(ii) وجود مجموعة الدوال كلها من~$\OLId{latin-looped}{A}$ إلى~$\OLId{latin-looped}{B}$.
\end{prob}

\begin{prob}
لتكن~$\OLId{latin-looped}{A}$ مجموعة، و~$\sim$ علاقة تكافؤ على~$\OLId{latin-looped}{A}$. برهن وجود مجموعة فئات
التكافؤ تحت~$\sim$ على~$\OLId{latin-looped}{A}$، أي~$\equivclass{\OLId{latin-looped}{A}}{\sim}$.
\end{prob}

\end{document}



